% Options for packages loaded elsewhere
\PassOptionsToPackage{unicode}{hyperref}
\PassOptionsToPackage{hyphens}{url}
\PassOptionsToPackage{dvipsnames,svgnames,x11names}{xcolor}
%
\documentclass[
  11pt,
  a4paper,
]{ltjsarticle}
\usepackage{amsmath,amssymb}
\usepackage{iftex}
\ifPDFTeX
  \usepackage[T1]{fontenc}
  \usepackage[utf8]{inputenc}
  \usepackage{textcomp} % provide euro and other symbols
\else % if luatex or xetex
  \usepackage{unicode-math} % this also loads fontspec
  \defaultfontfeatures{Scale=MatchLowercase}
  \defaultfontfeatures[\rmfamily]{Ligatures=TeX,Scale=1}
\fi
\usepackage{lmodern}
\ifPDFTeX\else
  % xetex/luatex font selection
\fi
% Use upquote if available, for straight quotes in verbatim environments
\IfFileExists{upquote.sty}{\usepackage{upquote}}{}
\IfFileExists{microtype.sty}{% use microtype if available
  \usepackage[]{microtype}
  \UseMicrotypeSet[protrusion]{basicmath} % disable protrusion for tt fonts
}{}
\makeatletter
\@ifundefined{KOMAClassName}{% if non-KOMA class
  \IfFileExists{parskip.sty}{%
    \usepackage{parskip}
  }{% else
    \setlength{\parindent}{0pt}
    \setlength{\parskip}{6pt plus 2pt minus 1pt}}
}{% if KOMA class
  \KOMAoptions{parskip=half}}
\makeatother
\usepackage{xcolor}
\usepackage[margin=24mm]{geometry}
\setlength{\emergencystretch}{3em} % prevent overfull lines
\providecommand{\tightlist}{%
  \setlength{\itemsep}{0pt}\setlength{\parskip}{0pt}}
\setcounter{secnumdepth}{5}
\ifLuaTeX
  \usepackage{selnolig}  % disable illegal ligatures
\fi
\IfFileExists{bookmark.sty}{\usepackage{bookmark}}{\usepackage{hyperref}}
\IfFileExists{xurl.sty}{\usepackage{xurl}}{} % add URL line breaks if available
\urlstyle{same}
\hypersetup{
  colorlinks=true,
  linkcolor={blue},
  filecolor={Maroon},
  citecolor={Blue},
  urlcolor={blue},
  pdfcreator={LaTeX via pandoc}}

\author{}
\date{}

\begin{document}

\hypertarget{ux6b63ux898fux65b9ux7a0bux5f0f}{%
\section{正規方程式}\label{ux6b63ux898fux65b9ux7a0bux5f0f}}

最小二乗問題

\[
f(x)=\|Ax-b\|_2^2
\]

を展開すると

\[
f(x)=x^\top A^\top Ax-2b^\top Ax+b^\top b
\]

です。

微分して

\[
\nabla f(x)=2A^\top Ax-2A^\top b
\]

を0と置くと

\[
\boxed{A^\top Ax=A^\top b}
\]

を得ます。

\hypertarget{ux5e7eux4f55ux5b66ux7684ux5c0eux51faux3068ux306eux4e00ux81f4}{%
\subsection{幾何学的導出との一致}\label{ux5e7eux4f55ux5b66ux7684ux5c0eux51faux3068ux306eux4e00ux81f4}}

これは残差が列空間に直交する条件

\[
A^\top(b-Ax)=0
\]

と同じです。

したがって微分による導出と射影による導出は、同じ事実を解析と幾何の二つの言語で表しています。

\hypertarget{ux6570ux5024ux8a08ux7b97ux4e0aux306eux6ce8ux610f}{%
\subsection{数値計算上の注意}\label{ux6570ux5024ux8a08ux7b97ux4e0aux306eux6ce8ux610f}}

正規方程式では \texttt{A\^{}TA}
を作るため、条件数がおおよそ二乗されます。

\[
\kappa_2(A^\top A)=\kappa_2(A)^2
\]

そのため数値的には QR や SVD の方が安全な場合があります。

\end{document}
